\documentclass[11pt]{article}
\usepackage[a4paper,margin=2.6cm]{geometry}
\usepackage{amsmath,amssymb,amsthm}
\usepackage[T1]{fontenc}
\usepackage{lmodern}
\usepackage{microtype}
\usepackage[hidelinks]{hyperref}
\newtheorem{theorem}{Theorem}
\newtheorem{lemma}[theorem]{Lemma}
\newtheorem{corollary}[theorem]{Corollary}
\theoremstyle{remark}
\newtheorem{remark}[theorem]{Remark}
\newcommand{\N}{\mathbb{N}_0}
\newcommand{\ud}{\bar d}
\newcommand{\ld}{\delta}
\title{Binary automata cannot certify a divergent $5n+p$ orbit}
\author{Research note, 25 November 2026}
\date{}
\begin{document}
\maketitle

\noindent\textbf{Status.} This note resolves none of the $155$ open census
cases. It proves that the certificate class used in the census can never
resolve one of them unless plain simulation already does. The argument has been
checked by its author and by the finite tests in the accompanying
verification script, and by nobody else. It uses
Konieczny's Theorem~B \cite{K}. Novelty has not been checked.

\section{Statement}

Fix an integer $p\ge1$ coprime to $10$ and let
\[
T(n)=n/2\ \ (n\text{ even}),\qquad T(n)=5n+p\ \ (n\text{ odd}).
\]
For a seed $s\ge1$, a \emph{first-step separator} is a set $S$ with
$T(S)\subseteq S$, $s\notin S$ and $T(s)\in S$; it proves $T^j(s)\ne s$ for all
$j\ge1$. A set is \emph{binary regular} if a finite automaton recognises the
binary expansions of its members. The congruence note shows that for
$5\nmid s$ no union of residue classes is a separator.

\begin{theorem}\label{main}
Let $5\nmid s$. A binary regular first-step separator for $s$ exists if and
only if the orbit of $s$ is eventually periodic and $s$ is not on its cycle.
\end{theorem}

When the orbit is eventually periodic the finite set $\{T^j(s):j\ge1\}$ is a
separator, and it is found by iterating. So for seeds coprime to $5$ a binary
automaton of any size proves nothing that simulation does not. In particular
each of the $155$ unresolved cases ($s=27$, $p$ prime) has a binary certificate
only if its orbit eventually cycles.

\begin{theorem}\label{structure}
Let $S$ be binary regular with $T(S)\subseteq S$. There is a periodic set $P$
such that $n\in P\iff T(n)\in P$ for every $n$ coprime to $5$, and such that
the integers coprime to $5$ on which $S$ and $P$ differ have density zero.
\end{theorem}

\section{Proofs}

Write $\ud$ for upper density, $g\simeq h$ when $g\ne h$ only on a set of upper
density zero, and $\ld$ for lower logarithmic density. $\ld$ is superadditive on
disjoint sets, at most $1$, and scales by $1/a$ under $n\mapsto an+c$. Call $n$
a \emph{unit} if $5\nmid n$. Units map to units. Let $f=\mathbf 1_S$.

\subsection{Backward words}

The preimages of $x$ are $2x$, and also $(x-p)/5$ when $x$ is even and
$x\equiv p\pmod 5$; the second preimage is odd. Write $d$ for $x\mapsto 2x$ and
$e$ for $x\mapsto(x-p)/5$. A \emph{word} is a finite sequence of these moves in
which every $e$ is immediately preceded by a $d$. A word $w$ with $J$ letters
$d$ and $o$ letters $e$ is defined on a set $C_w$ which is a union of residue
classes modulo $5^{o}$, apart from finitely many small $x$ where the value
would not be positive. There it is affine with slope $2^{J}/5^{o}$, and
$T^{|w|}(w(x))=x$. Different words of the same length have disjoint images,
because $w$ is the reversed parity itinerary of $w(x)$. Let $\mu$ denote the
density of a union of residue classes and put
\[
a_k=\sum_{|w|=k}5^{o(w)}2^{-J(w)}\mu(C_w),
\]
the sum over words that are empty or begin with $d$.

\begin{lemma}\label{mass}
$a_k=\tfrac23+\tfrac13(-\tfrac12)^k\ge\tfrac12$.
\end{lemma}
\begin{proof}
Such a word is $d^k$, contributing $2^{-k}$, or $d^{j}e\,w'$ with $j\ge1$ and
$w'$ empty or beginning with $d$. The move $d^je$ is defined on the class
$x\equiv p2^{-j}\pmod5$ and maps it bijectively onto the integers, so
$\mu(C_{d^jew'})=\mu(C_{w'})/5$. Its weight is $5\cdot2^{-j}$ times that of
$w'$. Hence $a_k=2^{-k}+\sum_{j=1}^{k-1}2^{-j}a_{k-1-j}$ with $a_0=1$, and the
closed form follows by induction.
\end{proof}

(As a check, $a_k$ is twice the density of $\{n:T^k(n)\text{ odd}\}$.)

\subsection{The entry set is thin on units}

Let $D=\{n:n\notin S,\ T(n)\in S\}$, a binary regular set. Since $f$ is
nondecreasing along orbits, the sets $T^{-k}(D)$, $k\ge0$, are pairwise
disjoint.

\begin{lemma}\label{equi}
Let $A$ be binary regular with $\ud(A)>0$. There are $P_0\ge1$, $a\ge0$ and
$c>0$ with $\ld(A\cap(r+5^k\N))\ge c\,5^{-k}$ for all $k\ge a$ and all
$r\equiv P_0\pmod{5^a}$.
\end{lemma}
\begin{proof}
This is the lemma on page 2 of the note \emph{Which binary regular sets are
forward closed?} with $3$ replaced by $5$. The proof only uses that appending a
binary word permutes the residues modulo the odd modulus: one finds an
accepting state $t$ in a closed component, two distinct loops $u,v$ of equal
length at $t$, and takes $a=v_5([u]_2-[v]_2)$, so that the loops generate the
translations by multiples of $5^a$ on the fibre over $t$.
\end{proof}

Suppose the set $A$ of units in $D$ has $\ud(A)>0$, and apply Lemma~\ref{equi}.
Shrinking to a subclass we may take $a\ge1$; then $P_0$ is a unit because $A$
consists of units. Since $2$ is a primitive root modulo $5$ there is $j\ge1$
with $2^jP_0\equiv p\pmod5$. Replacing $j$ by $j+4$ changes $(2^jP_0-p)/5$
modulo $5$ by the nonzero translation $3p$, so among five consecutive choices
at most one fails to give a unit. Repeating $a$ times gives a word
$w_0=d^{j_1}e\cdots d^{j_a}e$ whose domain is exactly the class of $P_0$ modulo
$5^a$, mapped bijectively onto the integers. For a word $w'$ that is empty or
begins with $d$ this gives $\mu(C_{w_0w'})=5^{-a}\mu(C_{w'})$. Each $C_{w_0w'}$
is a union of classes modulo $5^{o}$ inside the class of $P_0$, so
$\ld(A\cap C_{w_0w'})\ge c\,\mu(C_{w_0w'})$, and the image under the word has
$\ld$ at least $5^{o}2^{-J}$ times this. Summing over the disjoint images and
using Lemma~\ref{mass},
\[
\ld\bigl(T^{-k}(A)\bigr)\ \ge\ c\,2^{-J(w_0)}a_{k-|w_0|}\ \ge\ c\,2^{-J(w_0)-1}
\qquad(k\ge|w_0|).
\]
These sets are disjoint for different $k$, a contradiction. Hence
\begin{equation}\label{inv}
f(T(z))=f(z)\quad\text{for all units } z \text{ outside a set of upper density zero.}
\end{equation}

\subsection{Base change}

Let $F(n)=\star$ if $5\mid n$ and $F(n)=f(n)$ otherwise. This sequence over
$\{0,1,\star\}$ is $2$-automatic. Let $\mathcal K$ be the finite $2$-kernel of
$f$.

\begin{lemma}\label{bridge}
For $J\ge1$ and a unit $0<B<5^J$, the sequence $q\mapsto f(5^Jq+B)$ is $\simeq$
to $q\mapsto\kappa(5q+c)$ for some $\kappa\in\mathcal K$ and $0\le c<5$.
\end{lemma}
\begin{proof}
Put $m_0=5^Jq+B$ and $m_i=(2^{j_i}m_{i-1}-p)/5$ for $1\le i\le J-1$, choosing
$j_i\ge1$ so that the constant term $B_i$ of $m_i$ is a positive unit. This is
possible as above, since the constant term is determined modulo $5^{J-i}$ with
$J-i\ge1$. Then $m_{J-1}=5\cdot2^{E}q+B_{J-1}$ with $0<B_{J-1}<5\cdot2^{E}$,
every $m_i$ is an odd unit, and $T$ leads from $m_{J-1}$ back to $m_0$ through
the units $2^{l}m_{i}$. Every point on this path is an affine function of $q$
with positive slope, so \eqref{inv} gives $f(m_0)=f(m_{J-1})$ for almost all
$q$. Write $B_{J-1}=2^{E}c+R$ with $0\le R<2^{E}$ and
$\kappa(n)=f(2^{E}n+R)$.
\end{proof}

So every member of the $5$-kernel of $F$ is $\simeq$ to $F$, to the constant
$\star$, or to one of finitely many sequences $\kappa(5q+c)$. Thus $F$ is
$2$-automatic and asymptotically $5$-automatic, and Theorem~B of \cite{K} gives
a periodic sequence $\simeq F$. It equals $\star$ exactly on multiples of $5$.
Let $P$ be the set where it equals $1$. By \eqref{inv}, and because $T^{-1}$ of
a density-zero set has density zero, $\mathbf 1_P(T(z))=\mathbf 1_P(z)$ for
almost all units $z$. Both sides are periodic in $z$ on each parity class, so
the equality holds for all units. This proves Theorem~\ref{structure}.

\subsection{Proof of Theorem~\ref{main}}

If the orbit is eventually periodic and $s$ is not on the cycle, the finite set
$\{T^j(s):j\ge1\}$ is a separator. If $s$ is periodic there is no separator.
Suppose the orbit of $s$ is infinite, so that its terms tend to infinity, and
let $S$ be a binary regular separator with $P$ as above. A binary regular set
of density zero omits some block $U$: no member has $U$ as consecutive digits
after its leading $1$ (same argument as in the base-$2^x3^y$ note). Let $I_U$ be
the open arc of angles $\theta$ such that every real with
$\log_2$ congruent to $\theta$ modulo $1$ has binary digits $1U$ at the top.

\emph{Case $s\notin P$.} By invariance no $T^j(s)$ is in $P$, while all of them
are in $S$. So the orbit lies in the density-zero regular set of units in
$S\setminus P$. Let $o_1,o_2,\dots$ be its odd terms; there are infinitely many
and $o_{i+1}=(5o_i+p)2^{-\alpha_i}$. Hence
$\log_2o_{i+1}\equiv\log_2o_i+\log_25+\varepsilon_i\pmod1$ with
$\varepsilon_i\to0$. The set of limit points of $\log_2o_i$ modulo $1$ is
closed, nonempty and invariant under the irrational rotation by $\log_25$, so
it is the whole circle. Some $o_i$ has angle in $I_U$ and therefore contains
$U$. Contradiction.

\emph{Case $s\in P$.} The complement of $S$ is closed under taking preimages
and contains $s$. Let $\Gamma$ be the smallest set containing $s$ such that
$m\in\Gamma$ implies $2m\in\Gamma$ and, whenever $j\ge1$ and $(2^jm-p)/5$ is a
positive integer coprime to $5$, also $(2^jm-p)/5\in\Gamma$. Every member of
$\Gamma$ is a unit outside $S$, and lies in $P$ by invariance. So $\Gamma$ lies
in the density-zero regular set of units in $P\setminus S$. Choose
$y_0=s$ and $y_i=(2^{j_i}y_{i-1}-p)/5\in\Gamma$ with $j_1\ge J_0$ and
$j_i\ge3$; five consecutive admissible exponents are available at each stage
and at most one is excluded. Then
\[
y_g=\frac{2^{E_g}s}{5^{g}}\,(1-\eta_g),\qquad
0<\eta_g\le\frac ps\sum_{i\ge1}\frac{5^{i-1}}{2^{E_i}}<3p\,2^{-J_0},
\]
where $E_i=j_1+\dots+j_i$. The angles $\log_2s-g\log_25$ are dense modulo $1$.
Taking $J_0$ large and $g$ with angle well inside $I_U$ puts the angle of $y_g$
in $I_U$, so $y_g$ contains $U$. Contradiction. \qed

\section{Remarks}

\begin{remark}
The exact statement is available here, unlike in the base-$6$ note for the 1932
permutation, because in base $2$ every odd step turns $\log_2n$ by the same
angle $\log_25$ and halving does not turn it at all. Base $5$ and base $10$
have the same feature and the argument should adapt; this has not been written
out. Bases such as $20$ do not have it.
\end{remark}

\begin{remark}
For $5\mid s$ there are periodic separators (congruence note). Seeds with
$\gcd(s,p)>1$ are covered by the theorem as long as $5\nmid s$; for $p=3$,
$s=27$ the question is the $5n+1$ orbit of $9$.
\end{remark}

\begin{remark}[What a proof would have to look like]
Along a divergent orbit the parity sequence is not eventually periodic, since
an eventually periodic parity sequence determines a rational $2$-adic point
whose orbit is bounded or leaves the integers. Any certificate must therefore
be a forward-closed set that is not congruence-defined and not finite-state in
base $2$, and that still excludes the seed. No example of such a set is known
for any orbit of $5n+p$ with $\gcd(p,5)=1$.
\end{remark}

\begin{thebibliography}{9}
\bibitem{K} J.~Konieczny, \emph{An asymptotic version of Cobham's theorem},
arXiv:2209.09588, Definition~1.1 and Theorem~B.
\end{thebibliography}
\end{document}
